\section{Stations, Reservations, and the Limits of the Directory}

When Willits became a center, it inherited the problem of distance. The 1935 directory places a station at Covelo beneath Saint Anthony and labels it by the adjacent ``Indian Reservation.'' A 1940 entry names Celestine Quinlan as pastor, Sylvan Murphy as administrator, and stations at Covelo, Sherwood, and Laytonville reservations. The 1941 edition again lists Covelo, Sherwood, and Laytonville. These entries show that Saint Anthony's administrative reach extended into several Native-community geographies.

They do not show what happened there. ``Station'' is a Catholic administrative classification, not the name of a Native polity or a record of consent, participation, attendance, sacramental reception, cultural negotiation, or resistance. ``Indian Reservation'' is the directory editor's broad period label. It can point researchers toward a location while flattening distinct communities, governments, histories, and relations to the Catholic Church.

The Sherwood Valley Band's own public history supplies one corrective. The Band identifies continuing Pomo homeland and describes the rancheria's 1909 federal establishment and its successor relation to the former Mendocino Indian Reservation. That account should control the identity and continuity of the community where it speaks. It does not, however, narrate Saint Anthony's ministry. Putting the tribal source beside the Catholic directory exposes an important asymmetry rather than closing it: one source describes homeland and government; the other describes where priests believed they had pastoral responsibility.

Covelo presents a related risk. It lies in Round Valley, home to the Round Valley Indian Tribes and to a history larger than the parish's mission entry. A Catholic station there can be dated from annual directories, but those lines cannot establish that the station served only Native people or represent how Native Catholics understood it. Nor should the current mission of Our Lady Queen of Peace be projected backward unchanged into every earlier use of ``Covelo.'' The explicit Queen of Peace title first appears in the directory sequence located for 1969.

The correct historical question is not whether a list proves a successful mission. It is what kind of relationship the list claims, what evidence could test that claim, and whose voice is absent. Sacramental registers might reveal baptisms, marriages, sponsors, mobility, and the names by which people identified themselves, but registers require privacy-aware and community-sensitive use. Correspondence might reveal clerical aims, while tribal archives and oral histories might preserve reception or opposition. None was available here in a form that could responsibly complete the story.

\claimcheck{``Reservation stations document Catholic conversion of Native communities''}{They document a Catholic pastoral geography, not its outcome. The directories name assignments and locations. They do not supply Native testimony, participation counts, sacramental registers, consent, or consequences. This account therefore reports the stations while refusing to turn administrative reach into a triumphal narrative.}

The station record nevertheless changes the meaning of Saint Anthony. Willits was never only a church for an incorporated town. Its clergy operated at the meeting point of settler transport corridors, federal reservation and rancheria systems, Native homelands, rural settlements, and Catholic jurisdiction. Any deeper future parish history must be undertaken with the communities represented too thinly in the existing ecclesiastical archive.
